% Tapis de Quillen, with Quillen's offprint and preprint (cote 162-5, typed part): archive cote 162-5, site page 13 (PDF page 14).
% typescript. Draft compared: none.
% AI-transcribed by Claude Opus 5.5 (claude-opus-5-5, adaptive thinking, effort high), Message Batches API, batch t162_full (msgbatch_01UAh3zyWPd454AeDkhnmQ33), 2026-11-13; not yet checked by eye.
\scanpage{162-5/13}
\mspage{12}

X. De plus, on a la formule\inserted{\unsure{du type $g'f_{.} = f'_{.}g'$} ci-dessus, impliquant la formule} de projection habituelle. Ceci posé, les
constructions bien connues en théorie des correspondances permettent
de définir une catégorie $\underline{B}$ dont lds objets sont ceux de C, et les
Hom étant donnés par
\[ \mathrm{Hom}_{\underline{B}}(X,Y) = B(XxY) \quad . \]
On a un isomorphisme canonique $\underline{B} \simeq \underline{B}^{o}$ , qui est l'identité sur les
objets et la symmétrie $B(XxY)\simeq B(YxX)$ sur les flêches, et un foncteur
évident \xout{$\underline{E}$xxx$\underline{E}$x} $F_{.}$ : $C \to \underline{B}$ , qui est l'identité sur les objets et don
% NOTE: the right edge of the page is cut; the line ends with "don" (continuation not visible)
par les graphes sur les flêches. Composé avec la symmétrie précédente
on trouve un foncteur $F^{\cdot}$: \xout{Ex} $C \to \underline{B}$. On vérifie enfin que le couple
($F_{.}$,$F^{\cdot}$) satisfait aux conditions envisagées dans l'alinéa 1 (ceci aus
% NOTE: right edge cut; "aus" (continuation not visible)
devrait être formalisé au cas des correspondances génér ales ...). \unsure{On}
vérifie enfin qu'il a la propriété universelle qu'il faut. Lorsqu'on
se donne un couple ($F'_{!}$,$F'^{\cdot}$) à val eursdans une catégorie B', il faut
trouver une factorisation par B. Elle est toute trouvée sur les
objets, il faut la définir sur les flêches, donc définir, pour X,Y
dans C, une application
\[ (*) \qquad B(XxY) \to \mathrm{Hom}(X',Y') \quad . \]
Mais si Z est auddessus de XxY, d'où $p:Z \to X$ et $q:Z \to Y$, on lui assoc
% NOTE: right edge cut; "assoc" (continuation not visible)
un élément du deuxième membre de (*) comme le composé
\[ X' \xrightarrow{F'^{\cdot}(p)} Z' \xrightarrow{F'_{!}(\unsure{q})} Y' \quad . \]
Il faut prouver que l'élément obtenu en \unsure{remplaçant|remphçant} Z\hand{\unsure{,}} par un $Z_1$ cobord
% NOTE: right edge cut; "cobord" (continuation not visible); small handwritten marks ". 2. ," above "Z par"
dant au dessus de X\inserted{\unsure{ne varie pas|n'a pas varié}}. Or si $g$: $T \to XxY$ réalise le cobordisme, on
voit (regardant seulement T et la décomposition bord(T) = $Z \amalg Z_1$)
qu'il existe une variété $\overline{T}$ sur $S^1$ (obtenue en ``doublant'' T) \xout{etxxxx}
\xout{applicationxTxxxxx} trqnsversale au dessus d'un couple de poihts $s_o$,s
% NOTE: right edge cut; "$s_o$,s" (index of the second point not visible)
de $S^1$, tels que $Z_i$ soit l'image inverse \inserted{dans $\overline{T}$} de $s_i$ (i=0,1); d'autre part,

\margin{13}
